\section{Module Structure of \Frexlib{}}\label{app:first-appendix}
We summarises the implementation's code modules
and their relationship to this article:
\begin{center}
\begin{tabular}{ll}
  \begin{minipage}[t]{.44\textwidth}
\begin{itemize}[leftmargin=*]
\item[\textbullet] \AModule[ (\secref*{frexlib}--\subsecref{frex-via-coproducts})]{Frex}
  core definitions
\begin{itemize}[leftmargin=*]
\item[\textbullet] \AModule[ (\secref*{frexlib})]{Signature}
  operations \& arities
\item[\textbullet] \AModule[ (\secref*{frexlib},\subsecref{UPs})]{Algebra}
  algebraic structures and terms, homomorphisms
\item[\textbullet] \AModule[ (\secref*{frexlib})]{Presentation}
  axioms, equational theories
\item[\textbullet] \AModule[ (\secref*{frexlib})]{Axiom}
  common axiom schemes
\item[\textbullet] \AModule[ (\secref*{frexlib})]{Model}
  axiom-validating algebras
\item[\textbullet] \AModule[ (\subsecref{powers})]{Powers}
  parameterised algebras
\item[\textbullet] \AModule[ (\secref*{frexlib})]{Free}
  simplification in all algebras
\begin{itemize}[leftmargin=*]
\item[\textbullet] \AModule[ (\subsecref{UPs})]{Definition}
  universal property
\item[\textbullet] \AModule[ (\secref*{certification})]{Construction} a
  non-effective quotient construction used for extraction, printing,
  and certification
\begin{itemize}[leftmargin=*]
\item[\textbullet] \AModule[ (\subsecref{fral via frex})]{ByFrex}
  reuse a frex simplifier to define a fral simplifier
\item[\textbullet] \AModule[ (\subsecref{extracting
    certificates}--\subsecref{proof simplification})]{Linear}
  generic proof simplification and printing
\item[\textbullet] \AModule[ (\secref*{certification})]{Idris}
  generic certification
\end{itemize}
\end{itemize}
\item[\textbullet] \AModule[ (\subsecref{frex-via-coproducts})]{Coproduct}
  universal property
\item[\textbullet] \AModule[ (\subsecref{UPs}--\subsecref{frex-via-coproducts})]{Frex}
  universal property, reuse coproduct and fral simplifier to define a frex simplifier
\begin{itemize}[leftmargin=*]
\item[\textbullet] \AModule[ (\secref*{certification})]{Construction}
  non-effective quotient construction used for extraction, printing,
  and certification
\end{itemize}
\item[\textbullet] \AModule[ (\secref*{certification})]{Lemma}
  auxiliary representation for auxiliary lemmata discharged by fral
  simplifiers, printed, or certified
\item[\textbullet] \AModule[ (\secref*{reflection})]{Magic}
  generic reflection code for ergonomic invocation
\end{itemize}
\end{itemize}
\end{minipage}
&
\begin{minipage}[t]{.52\textwidth}
\begin{itemize}[leftmargin=*]
\item[\textbullet] \AModule{Frexlet.Monoid}
  modules concerning varieties of monoids and their simplifiers
\begin{itemize}[leftmargin=*]
\item[\textbullet] \AModule[ (\secref*{frexlib})]{Theory}
  signature, axioms, pretty printing for the theory of ordinary monoids
\item[\textbullet] \AModule[ (\secref*{frexlib})]{Notation}
  shared infix notation (additive and multiplicative) for monoid varieties
\item[\textbullet] \AModule[ (\figref{monoid frexlets})]{Frex}
  frex simplifier for monoids
\item[\textbullet] \AModule[ (\subsecref{fral via frex})]{Free}
  fral simplifier, reuses frex simplifier
 %  -- Concrete algebras
\item[\textbullet] \AModule[ (\secref*{frexlib})]{Nat}
  additive and multiplicative monoid structure of the natural numbers
\item[\textbullet] \AModule{Pair} types with the cartesian product as
  a proof-relevant monoid structure
\item[\textbullet] \AModule{List}
  monoid structure of lists with catenation
\item[\textbullet] \AModule{Commutative}
  commutative monoids modules
\begin{itemize}[leftmargin=*]
\item[\textbullet] \AModule{Theory}
  commutativity axiom
\item[\textbullet] \AModule[ (\subsecref{UPs}]{Free})
  fral simplifier
\item[\textbullet] \AModule[ (\subsecref{UPs})]{NatSemiLinear}
  auxiliary definitions for fral simplifier
\item[\textbullet] \AModule[ (\subsecref{frex-via-coproducts})]{Frex}
  simplifier, reuses fral via coproducts
\item[\textbullet] \AModule[ (\subsecref{frex-via-coproducts})]{Coproduct}
  coproduct of commutative monoids
\item[\textbullet] \AModule{Nat}
  addition and multiplication of naturals
\end{itemize}
\item[\textbullet] \AModule{Involutive}
  modules concerning monoids equipped with an involution
\begin{itemize}[leftmargin=*]
\item[\textbullet] \AModule[ (\secref*{frexlib})]{Theory}
  signature and axioms
\item[\textbullet] \AModule{Free}
  simplifier, reuses frex simplifier
\item[\textbullet] \AModule[ (\subsecref{reusing frexlets})]{Frex}
  simplifier, reuses monoid frex
\item[\textbullet] \AModule[ (\secref*{frexlib})]{List}
  involutive monoid structure of list reversal
\end{itemize}
\end{itemize}
\end{itemize}
\end{minipage}
\end{tabular}
\end{center}


\section{Extensional Function and Quotient Setoids}
\label{sec:setoid example}
\applabel{setoid example}
\begin{figure*}
  \TwoColDIYprime{.45}{
 \QuotientDef[numbers=right, numbersep=0.2cm]{}
 \FunSpaceDef[numbers=right, firstnumber=17,numbersep=0.2cm]{}
  }{.54}{%
    \VectEquality{}%
    \VectSetoid{}%
    \VectMap{}%
  }
  \caption{(a) Quotient, function-space, and (b) vector setoids (top) and a higher-order homomorphism (bottom)}
  \figlabel{quotients-and-functions}
  \figlabel{vect-setoid}
\end{figure*}
\Figref{quotients-and-functions}a defines the quotient of a type by a
function \QuotientingMap{}, taking two elements to be equal when their
images under the function \IdrisBound{q} are equal, and
the setoid of homomorphisms
between two setoids together with extensional equality.  This example
also demonstrates \idris's local definitions (lines 4--6, e.g.), possibly
with quantities, named-argument function calls (lines 8--15, e.g.), application operator \IdrisBound{\$},
and
anonymous functions (lines 9--10, e.g.). \idris, like \Haskell, implicitly
quantifies (with quantity \IdrisKeyword{0}) over unbound variables in
type-declarations such as the type \ImplicitArg{} in \QuotientName{}.
These underscores mean elaboration must fill-in the blanks
uniquely using unification.


\Figref{vect-setoid}b presents a setoid over $\nvar$-length
vectors over a given setoid. The vector functorial action
\VectMapName{} has a setoid homomorphism structure
between the two setoids of homomorphisms: (1) \VectMapF{} is a homomorphism (lines 19--25), and
that (2) it maps extensionally equal homomorphisms to extensionally
equal homomorphisms (26--30). These proofs use \idris's equational
reasoning notation for setoids (lines 20--25 and 27--30), a
deeply embedded chain of equational steps. Each step
\IdrisData{(\KatlaTilde\KatlaTilde)} appeals to
transitivity, and requires a justification. The last two dots in the
thought bubble operator \BasicThoughtBubble{} modify the reason:
plain usage (line~23) appeals to a setoid equivalence;
an
equals in the middle dot, e.g.~\EqualThoughtBubble{}, appeals to reflexivity via propositional equality (lines 22, 25, 28, 30); and
a comparison symbol in the end, e.g.~\EqualSymThoughtBubble{}, appeals to
symmetry (lines 25, 30).


\section{Proof Printing and Certification}
\applabel{last-appendix}
\applabel{printing}
\Figref{linear-derivations} presents the layered representation of
linear derivations.
%
\begin{figure}
\TwoColDIY{.4}{%
  \RTList
  \caption*{(a) reflexive-transitive closure}
  \Symmetrise
  \caption*{(b) symmetric closure}
  \vspace{2\baselineskip}
  \Derivation%
  \caption*{(e) linear derivations}
}{.57}{%
  \Locate
  \caption*{(c) unary congruence closure}
  \Step
  \caption*{(d) axiomatic steps}
}
\caption{Layered (a--d) representation of linear derivations (e)}
\figlabel{linear-derivations}
\end{figure}

\figref{proof-extraction} shows an automatically extracted proof for
the equation $(x•3)•2 = 5•x$ in the additive monoid structure $(\Nat,
0, (\Plus))$. The extracted proof has $\IntroPrinterStepCount$ steps ---
 far from the shortest proof possible.
Extraction removes reflexivity and transitivity steps, and the pointed
bracket tells whether the step uses the axiom directly (angle points
right) or using symmetry (angle points left). Square brackets mean
appealing to congruence, where the context is the congruence's
context, and the term in the hole is the equation's LHS.
\begin{figure*}
  \IntroPrintedListing{}
  \caption{\Frexlib-extracted proof of $(x•3)•2 = 5•x$
    in the additive monoid over \Nat{}}
  \figlabel{proof-extraction}
\end{figure*}
\figref{certification} shows an automatically extracted certificate
for the equation $0 + (x + 0) + 0 = x$ in a generic monoid $\mvar{} =
(\Um{}, \OI{}, (\DotPlus))$.  The certificate is generated inside a
module that parameterises over the generic monoid \mvar{} and
introduces the various notations and reasoning functions.

\begin{figure*}
  \Certificate{}
  \caption{\Frexlib-certificate for the  of $0 + (x + 0) + 0 = x$ in a generic monoid \mvar{}}
  \figlabel{certification}
\end{figure*}
